%% ================================================================
%% AX5 Non-Reducibility Remark
%% To be inserted immediately after \begin{definition}[DSTP] in
%% paper1.tex, i.e., after Definition~\ref{def:dstp}.
%% ================================================================

\begin{remark}[Non-reducibility of \textup{AX5} — constructive witness]
\label{rem:ax5_independence}
\textup{AX5} is logically independent of \textup{AX1--AX4}.
The independence is witnessed constructively as follows.

\smallskip
\noindent\textbf{Witness rule.}
Let $\{p_j = -T + 2Tj/N\}_{j=1}^N$ be the equidistant Riemann rule
with weights $w_j = 2T/N$.
This rule satisfies:
\begin{itemize}
  \item \textup{AX3} (first-order quadrature consistency), with
    $C_Q = 1$ and $h = 2T/N$;
  \item \textup{AX4} (uniform PSWF bounds \cite{BonamiKaroui2014}),
    since $\|\psi_n^{(c)}\|_{L^\infty}$ and
    $\|(\psi_n^{(c)})'\|_{L^\infty}$ are intrinsic to the PSWFs and
    independent of the quadrature nodes.
\end{itemize}

\noindent\textbf{AX5 fails.}
For this rule,
\[
  \sup_{m \ne n,\; m,n \le N-1} |E_{mn}| \gtrsim hc,
\]
with no decay in $|m-n|$: the equidistant nodes carry no
phase-alignment with the oscillatory structure of $\psi_m^{(c)}
\overline{\psi_n^{(c)}}$, so the quadrature error
$R_{mn}^{\mathrm{quad}}$ cannot exploit destructive interference.
Consequently, $\|E_{N,c}\|_2 \lesssim hc \cdot N$ (Schur test with
no off-diagonal decay), and for $N \sim c$ this diverges,
yielding no frame stability.

\smallskip
\noindent\textbf{Structural interpretation.}
The failure of \textup{AX5} for the Riemann rule is not a
consequence of low approximation order or insufficient regularity.
\textup{AX1--AX4} together supply:
\begin{enumerate}[(i)]
  \item $c$-uniform symbol-class control of the PSWF amplitude and
    phase (AX1, proved in \cite{PaperIV});
  \item spectral gap $|\chi_m - \chi_n| \ge \kappa_1\max(1,c)|m-n|$
    (AX2, algebraic consequence of \cite{OsipovRokhlinXiao});
  \item first-order quadrature consistency (AX3);
  \item uniform $L^\infty$ and derivative bounds on PSWFs (AX4).
\end{enumerate}
None of these conditions encodes \emph{geometric compatibility}
between the node set $\{p_j\}$ and the oscillatory eigenstructure
of $V_{N,c}$. \textup{AX5} is not an estimate of higher order;
it is a condition of a different type entirely:

\begin{quotation}
  \noindent\textit{AX5 is not an estimate. It is a geometric
  compatibility condition between oscillatory eigenstructure and
  discrete sampling.}
\end{quotation}

This places \textup{AX5} in the same logical category as spectral
exactness conditions in Gauss quadrature, sampling compatibility
conditions in the Shannon--Nyquist theorem, and frame reconstruction
axioms in non-uniform sampling theory --- all of which assert
structural resonance between a discrete node set and a continuous
function class, and none of which follow from regularity bounds alone.

\smallskip
\noindent\textbf{Consequence for the axiom system.}
\textup{AX5} is not ``stronger'' than \textup{AX1--AX4}; it is
\emph{orthogonal} to them in the following precise sense:
\textup{AX1--AX4} control the \emph{analytic} structure of the
system (regularity, spectral gaps, approximation order);
\textup{AX5} controls the \emph{geometric} compatibility of the
discretisation. No combination of analytic bounds implies geometric
compatibility. Hence the axiom system $\{\textup{AX1},\dots,
\textup{AX5}\}$ is minimal: no axiom is derivable from the others.

Verification of \textup{AX5} for PSWF-adapted (XRY/Gauss--PSWF)
quadrature schemes is the content of~\cite{PaperIIquad}.
\end{remark}
